\chapter{Reto final integrador}
\sourcebadge{Reconstruido desde el índice del reto integrador.}

El reto conecta los elementos del caso en una cadena verificable:

\begin{center}
\textbf{Problema} $\rightarrow$ \textbf{Usuario y tarea} $\rightarrow$ \textbf{Evento} $\rightarrow$ \textbf{Característica ISO} $\rightarrow$ \textbf{Métrica} $\rightarrow$ \textbf{Resultado} $\rightarrow$ \textbf{Riesgo} $\rightarrow$ \textbf{Recomendación}
\end{center}

Para la facturación duplicada, el usuario es admisión, la tarea es emitir una factura, el evento es \texttt{SIM-BILL-DUP}, la característica es libertad de riesgo y la métrica M-RIE-01 resulta 4,50\% frente a una meta menor al 1\%. El riesgo es financiero y la recomendación es probar idempotencia y conciliación en una iteración posterior, sin alterar la línea base.

La misma estructura se aplica a HCE, citas, prescripción, imagenología y telemedicina. Esto permite defender cada conclusión con una fuente concreta y evita recomendaciones basadas únicamente en percepción.

